\subsection{The leading Cartier term of a possible point kernel}
\paragraph{A2V (Proved necessary condition; not a solution of A2).}
Let $X/k$ be smooth, projective and integral over an algebraically
closed field of characteristic $p>0$, with $\dim X\geq2$.
Let $Y\subset X$ be a smooth ample divisor and put
$L=\mathcal O_Y(Y)$. Suppose a nonzero closed global one-form
$\eta$ vanishes along $Y$. Let $m\geq1$ be its exact vanishing
order along the integral divisor $Y$. Then $p\mid m$.
If in addition $C\eta=\eta$, write $m=pa$. There is a nonzero
element
\[
 H^0(Y,B_Y^1\otimes L^{-a})
  \simeq
 \ker\bigl(F:H^1(Y,L^{-a})\longrightarrow
                   H^1(Y,L^{-pa})\bigr),                 \tag{A2V.1}
\]
where $B_Y^1=\operatorname{im}(d:F_*\mathcal O_Y\to
F_*\Omega_Y^1)$ has its Frobenius-module structure.
The cohomology group in the Frobenius target is understood
with the corresponding Frobenius scalar structure; equivalently,
the map is $p$-semilinear on the usual $k$-vector spaces.
In particular a nonzero $p$-torsion line bundle on $X$ that
restricts trivially to $Y$ forces this obstruction for some $a\geq1$.

Here is the proof. The leading coefficient of $\eta$
is a nonzero section
\[
 \sigma\in H^0\bigl(Y,\Omega_X^1|_Y\otimes L^{-m}\bigr).
\]
Locally write $\eta=t^m\alpha$, where $t$ cuts out $Y$.
The equation $d\eta=0$ implies
$m\,dt\wedge\alpha=0$ modulo $t$. If $p\nmid m$, the
leading coefficient is purely conormal. The conormal sequence
would put it in
\[
 H^0(Y,L^{-m-1})=0,
\]
contrary to its definition. Thus $m=pa$.

Since $d(t^{pa})=0$, the local form $\alpha$ is closed.
Cartier's rule $C(f^p\alpha)=fC(\alpha)$ and $C\eta=\eta$
give
\[
 C\alpha=t^{(p-1)a}\alpha.
\]
Restriction to the smooth divisor commutes with Cartier, so
the tangential leading coefficient $\beta$ has $C_Y\beta=0$.
It is nonzero: otherwise $\sigma$ would again be a section of
$L^{-pa-1}$, which has no nonzero global section. When local
equations change by $t_i=u_{ij}t_j$, the coefficients change
by $\alpha_i=u_{ij}^{-pa}\alpha_j$. Because the multiplier is
a $p$th power, closedness and exactness are preserved, and the
Frobenius-module interpretation makes these precisely the
transition functions of $B_Y^1\otimes L^{-a}$. The exact
Cartier sequence identifies $\ker C_Y$ on closed forms with
$B_Y^1$.

Finally, tensoring the Frobenius sequence with $L^{-a}$ gives
\[
 0\longrightarrow L^{-a}\longrightarrow F_*L^{-pa}
  \longrightarrow B_Y^1\otimes L^{-a}\longrightarrow0.
\]
Both negative line bundles have no global section, so its
cohomology sequence proves (A2V.1). The two exact sequences
and projection-formula conventions can also be checked in
\cite[Lemmas 9.2 and 9.5, pp.~2106--2107]{DGJ_Langer2011}.
For a point kernel take $\eta=d\log(M)$. It is nonzero and
Cartier-fixed; triviality of $M|_Y$ first makes its tangential
restriction zero, and the conormal sequence with
$H^0(Y,L^{-1})=0$ makes its full restriction zero.

There is also a numerical exclusion that can be checked on the
ambient variety. Since $m\geq p$, this same nonzero form defines
an injection of sheaves
\[
 \mathcal O_X(pY)\hookrightarrow\Omega_X^1.
\]
It is injective because $X$ is integral and the map is nonzero
at its generic point. Taking global sections gives
\[
 h^0(X,\Omega_X^1)\ \geq\ h^0(X,\mathcal O_X(pY)).
                                                        \tag{A2V.2}
\]
Consequently, if $\dim X\geq3$, the reverse strict inequality
excludes every nontrivial point in the restriction kernel by
the already established finite-$p$-group reduction. This bound
does not assert that such an inequality holds for all $X$.

\paragraph{Why Cartier iteration alone stops here.}
For every prime $p$, the affine local model
\[
 U=\operatorname{Spec}k[t,x,(1+t^px)^{-1}],\qquad
 \eta=d\log(1+t^px)=\frac{t^p\,dx}{1+t^px}
\]
is closed and Cartier-fixed and has exact order $p$ along
$t=0$. Its tangential leading term is the nonzero exact form
$dx$. Explicitly,
\[
 \eta=\sum_{n\geq0}(-1)^n t^{p(n+1)}x^n\,dx.
\]
Cartier retains $n=pr+p-1$ and transforms those terms into
$(-1)^r t^{p(r+1)}x^r\,dx$, giving the same series.
Thus one cannot infer divisibility of the vanishing order by
$p^2$, or infinite-order vanishing, from Cartier fixedness
alone. This is an affine local model, not a projective ample
counterexample. Its exact finite coefficient identities for
$p=2,3,5,7$ are recorded in
\texttt{check\_cartier\_leading.py} and the saved JSON output.

\paragraph{A necessary singularity of the associated cover.}
Any actual point-kernel counterexample has a nonnormal
$\mu_p$-torsor along $Y$. Indeed, the torsor has local equation
$z^p=u$, with $u$ a unit and $du/u=\eta$. It is integral:
if $u$ were a $p$th power in the function field, the generic
torsor section would extend over the normal scheme $X$,
trivializing $M$. Since $\eta$ vanishes along $Y$, the
Jacobian criterion makes this hypersurface singular at the
generic point over $Y$. That is a codimension-one point.
The integral hypersurface is Cohen--Macaulay, but fails
regularity in codimension one, so it is not normal. In
particular constructions whose $\mu_p$-cover stays normal
cannot yield the required point kernel. This observation
does not rule out the nonnormal-cover route.

\paragraph{Exact remaining gate.}
The necessary Frobenius kernel on the smooth ample divisor
need not vanish in positive characteristic. Its existence
does not construct a Cartier-fixed form on an ambient smooth
projective variety, and no converse is proved here. The
remaining task is either to exclude that global extension or
to realize it with a fully specified smooth very ample pair.
